\documentclass[../../../include/open-logic-section]{subfiles}

\begin{document}

\olfileid{sth}{ordinals}{vn}
\olsection[વૉન નોયમનની રચના]{વૉન નોયમનની ક્રમવાચક સંખ્યાઓની રચના}

\olref[sth][ordinals][iso]{thm:woalwayscomparable}થી એક વિચાર
સૂઝે છે. સુક્રમોના અવેજ તરીકે કામ કરવા માટે
\emph{ક્રમપ્રકારો} કહેવાતી અમુક વસ્તુઓ રજૂ કરી શકીએ.
સુક્રમ $\tuple{A, <}$ના ક્રમપ્રકાર માટે
$\ordtype{A, <}$ લખીએ, તો નીચેના બે સિદ્ધાંતો મળે
એવી આશા રાખીએ:
\begin{align*}
	\ordtype{A, <} = \ordtype{B, \lessdot} &
	\text{ તો અને ફક્ત તો } \ordeq{\tuple{A, <}}{\tuple{B, \lessdot}}\\
	\ordtype{A, <} < \ordtype{B, \lessdot}&
	\text{ તો અને ફક્ત તો }\ordeq{\tuple{A, <}}{\tuple{B_b, \lessdot_b}}\text{ જ્યાં કોઈ }b \in B
\end{align*}
વધુમાં, જેમ પ્રાકૃતિક સંખ્યાઓને અમુક ગણો તરીકે રજૂ
કરી શકીએ છીએ, તેમ ક્રમપ્રકારોને પણ \emph{અમુક ગણો
તરીકે} રજૂ કરી શકીએ એવી આશા રાખીએ.

આવું કરવાની સૌથી સામાન્ય રીત---અને આપણે અનુસરવાની
રીત---એ આ ક્રમપ્રકારોને અમુક \emph{પ્રમાણભૂત} સુક્રમિત
ગણો દ્વારા વ્યાખ્યાયિત કરવાની છે. આ પ્રમાણભૂત ગણો
સૌપ્રથમ વૉન નોયમને રજૂ કર્યા:

\begin{defn}
ગણ $A$ \emph{પરંપરિત} હોય જો અને ફક્ત જો
$(\forall x \in A)x \subseteq
A$. પછી $A$ એ \emph{ક્રમવાચક સંખ્યા} હોય જો અને ફક્ત
જો $A$ પરંપરિત હોય અને $\in$ વડે સુક્રમિત હોય.
\end{defn}
\noindent
આગળ ક્રમવાચક સંખ્યાઓ માટે ગ્રીક અક્ષરો વાપરીશું.
વ્યાખ્યાથી સીધું મળે છે કે જો $\alpha$ ક્રમવાચક
સંખ્યા હોય, તો $\tuple{\alpha, \in_\alpha}$
સુક્રમ છે, જ્યાં $\in_\alpha =
\Setabs{\tuple{x, y} \in \alpha^2}{x \in y}$.
તેથી સંકેતનો થોડો છૂટો ઉપયોગ કરીને આપણે કહી
શકીએ કે $\alpha$ \emph{પોતે} સુક્રમ છે.

આપણી શરૂઆતની કેટલીક ક્રમવાચક સંખ્યાઓ આ છે:
\[
	\emptyset, \{\emptyset\},
	\{\emptyset, \{\emptyset\}\},
	\{\emptyset, \{\emptyset\}, \{\emptyset, \{\emptyset\}\}\}, \ldots
\]
તમે નોંધશો કે આ જ કેટલીક શરૂઆતની ક્રમવાચક
સંખ્યાઓ અનંતતાની સ્વયંસિદ્ધિમાં, એટલે કે વૉન
નોયમનની $\omega$ની વ્યાખ્યામાં, મળી હતી
(\olref[sth][z][infinity-again]{sec} જુઓ).
આ કોઈ યોગાનુયોગ નથી. વૉન નોયમનની ક્રમવાચક
સંખ્યાઓની વ્યાખ્યા પ્રાકૃતિક સંખ્યાઓને ક્રમવાચક
સંખ્યાઓ તરીકે લે છે, પણ અનંતાતીત ક્રમવાચક
સંખ્યાઓને પણ સ્થાન આપે છે.

હંમેશની જેમ હવે પૂછી શકીએ: શું આ જ ક્રમવાચક
સંખ્યાઓ \emph{છે}? કે વૉન નોયમને ફક્ત એવા
કેટલાક ગણો આપ્યા છે જેને આપણે ક્રમવાચક
સંખ્યાઓ \emph{તરીકે લઈ} શકીએ?
આ પ્રશ્ન અંગેની ચર્ચા
\olref[sfr][rel][ref]{sec},
\olref[sfr][arith][ref]{sec},
\olref[sfr][infinite][dedekindsproof]{sec} અને
\olref[sth][z][nat]{sec}માં થયેલી ચર્ચાઓ જેવી
જ હશે; તેથી ફરી વિસ્તાર નહીં કરીએ.
આગળ “ક્રમવાચક સંખ્યાઓ”નો અર્થ માત્ર “વૉન
નોયમનની ક્રમવાચક સંખ્યાઓ” જ લઈશું.

\end{document}
